\section{Tree Ensemble Extensions}\label{app:extensions}

This appendix records how the proof skeleton adapts to nearby tree-ensemble
trainers. In both extensions, the forest encoding, leaf-routing lookups,
well-formedness checks, histogram construction, and the two batching techniques
remain unchanged. What changes is the plaintext trainer: the training target,
the split objective, the leaf-value formula, and the round update.

\subsection{Extending to XGBoost}\label{app:xgboost-extension}

\paragraph{Statement.}
The public statement fixes an XGBoost-style~\cite{xgboost16} quantized trainer,
including the loss function, learning rate $\eta$, regularization parameters
$\lambda$ and $\gamma$, tree-shape policy, quantization parameters, and all
dimensions.  The prover commits to the dataset, labels, forest, leaf values,
and the training trace.  The verifier should be convinced that the committed
forest is obtained by running this public XGBoost-style trainer on the
committed data.

\paragraph{Plaintext trainer.}
XGBoost also trains an additive model over rounds, but the tree for each class
is fitted using a second-order approximation of the loss.  As in
\Cref{eq:plain-score-update}, write $\mathbf{S}[t-1,k,r]$ for the score of sample $r$
and class $k$ before the tree of round $t$ is added, with
$\mathbf{S}[-1,k,r]=0$.  The public trainer fixes the loss and therefore fixes the
sample-wise gradient and Hessian values
\[
    g_{t,k,r}
    =
    \frac{\partial \mathcal{L}^{(t-1)}}{\partial \mathbf{S}[t-1,k,r]},
    \qquad
    h_{t,k,r}
    =
    \frac{\partial^2 \mathcal{L}^{(t-1)}}{\partial \mathbf{S}[t-1,k,r]^2}.
\]
For multi-class classification, we use the standard class-wise diagonal-Hessian
approximation, so each class-$k$ tree in round $t$ is trained from the vector
$(g_{t,k,r},h_{t,k,r})_{r\in[M]}$.

Fix a node $u$ in the class-$k$ tree of round $t$, and let $I_u$ be the samples
reaching it.  For a candidate split specified by feature $j$ and threshold bin
$b$, define
\begin{align}
L_{u,j,b} &:= \{\, r\in I_u : \mathbf{D}[r,j] < b \,\},\\
R_{u,j,b} &:= \{\, r\in I_u : \mathbf{D}[r,j] \ge b \,\}.
\end{align}
For any sample set $A\subseteq[M]$, write
\[
    G_A := \sum_{r\in A} g_{t,k,r},
    \qquad
    H_A := \sum_{r\in A} h_{t,k,r}.
\]
Define $G_u=G_{I_u}$, $H_u=H_{I_u}$,
$G_L=G_{L_{u,j,b}}$, $H_L=H_{L_{u,j,b}}$,
$G_R=G_{R_{u,j,b}}$, and $H_R=H_{R_{u,j,b}}$.
The proof uses the following score, which is twice the usual XGBoost gain and
therefore gives the same ordering and pruning decision:
\begin{equation}\label{eq:xgb-gain}
    \Phi_{\mathsf{xgb}}(u,j,b)
    =
    \frac{G_L^2}{H_L+\lambda}
    +
    \frac{G_R^2}{H_R+\lambda}
    -
    \frac{G_u^2}{H_u+\lambda}
    -
    2\gamma .
\end{equation}
The trainer selects the valid candidate with largest positive score; if no such candidate exists, the node becomes a
leaf.  Once the tree structure is fixed, the leaf value of a leaf $\ell$ with
sample set $I_\ell$ is
\begin{equation}\label{eq:xgb-leaf}
    w_{\ell}
    =
    -
    \frac{G_{I_\ell}}{H_{I_\ell}+\lambda}.
\end{equation}
The score update is
\[
    \mathbf{S}[t,k,r]
    =
    \mathbf{S}[t-1,k,r]
    +
    \eta \cdot w_{\ell^{(t,k)}(r)} ,
\]
where $\ell^{(t,k)}(r)$ is the leaf reached by sample $r$ in the tree trained
at round $t$ for class $k$.

\paragraph{Proof sketch.}

For multi-class logistic loss, the gradient and Hessian in round $t$ are
computed from the current prediction scores, which depend on all trees trained
in the previous rounds. More precisely, let $\mathbf{S}[t-1,k,r]$ denote the
score of sample $r$ for class $k$ before round $t$ starts. It is the accumulated
output of the class-$k$ trees from all earlier rounds:
\[
    \mathbf{S}[t-1,k,r]
    =
    \sum_{\tau < t}
    \eta \cdot
    w_{\ell^{(\tau,k)}(r)},
\]
where $\ell^{(\tau,k)}(r)$ is the leaf reached by sample $r$ in the class-$k$
tree trained at round $\tau$, and $w_{\ell^{(\tau,k)}(r)}$ is the corresponding
leaf value. We set $\mathbf{S}[-1,k,r]=0$ if the rounds are indexed from $0$.

The plaintext trainer first applies softmax to these accumulated scores:
\[
    p_{t,k,r}
    =
    \frac{\exp(\mathbf{S}[t-1,k,r])}
    {\sum_{c\in[K]}\exp(\mathbf{S}[t-1,c,r])}.
\]
Under the standard diagonal-Hessian approximation used by
XGBoost~\cite{xgboost16}, the sample-wise gradient and Hessian for class $k$
are
\[
    g_{t,k,r}=p_{t,k,r}-\mathbf{Y}[k,r],
    \qquad
    h_{t,k,r}=p_{t,k,r}(1-p_{t,k,r}),
\]
where $\mathbf{Y}$ is the one-hot label matrix. These values are then
aggregated over the samples reaching each node and used to compute the split
gain and leaf values.

Therefore, to prove the correctness of $\mathbf{G}$ and $\mathbf{H}$, the
prover first proves that the committed score tensor $\mathbf{S}$ is obtained by
correctly accumulating the leaf values of all previous trees along the committed
routing paths. Given $\mathbf{S}$, the prover proves that the committed
probability tensor $\mathbf{P}$ is the softmax of $\mathbf{S}$ using the same
lookup and fixed-point gadgets as in the main construction. Finally, it proves
the pointwise identities
\begin{align}
    \mathbf{G}[t,k,r]
    &=
    \mathbf{P}[t,k,r]-\mathbf{Y}[k,r],\\
    \mathbf{H}[t,k,r]
    &=
    \mathbf{P}[t,k,r]\cdot(1-\mathbf{P}[t,k,r]),
\end{align}
with the required fixed-point rescaling for the multiplication. The linear
relations are batched by domain-lifting batching, while the multiplication,
lookup, and quantization checks are batched by interleaving batching.

After this step, the rest of the proof follows the same structure as the GBDT
proof: the prover proves the histogram transformation from sample-wise
gradients and Hessians to per-node frequency statistics, propagates the
histograms to internal nodes, proves the XGBoost gain for each candidate split,
and proves that the selected split is optimal.

\subsection{Extending to Random Forests}\label{app:rf-extension}

\paragraph{Statement.}
The public statement fixes a quantized random-forest trainer, including the
number of trees, tree-shape policy, split objective, leaf-value rule,
quantization parameters, and the randomness used for sampling data points and
feature subsets.  The randomness can either be public directly or derived from a
public seed.  The prover commits to the dataset, labels, forest, leaf values,
and the proof-facing training trace.  The verifier should be convinced that the
committed forest is obtained by running this public random-forest trainer on the
committed data.

\paragraph{Plaintext trainer.}
Unlike GBDT and XGBoost, random forests have no boosting state.  Trees are
trained independently.  For each tree $q$, the public trainer fixes a sample
multiplicity vector
\[
    \boldsymbol{\rho}^{(q)}\in\mathbb{F}^{M},
\]
where $\boldsymbol{\rho}^{(q)}[r]$ records how many times sample $r$ is included
in the bootstrap sample, and a feature-availability mask
\[
    \mathbf{F}^{(q)}\in\{0,1\}^{N_{\mathrm{tot}}\times d},
\]
where $\mathbf{F}^{(q)}[u,j]=1$ means that feature $j$ is allowed when splitting
node $u$ of tree $q$.  If the implementation uses sampling without replacement,
then $\boldsymbol{\rho}^{(q)}$ is simply a boolean mask.

For classification, the natural histogram object is a class-frequency
histogram.  For each node $u$, feature $j$, bin $b$, and class $c$, define
\[
    \mathbf{N}[u,j,b,c]
    =
    \sum_{r\in[M]}
    \boldsymbol{\rho}^{(q)}[r]\cdot
    \mathbf{1}\{r\text{ reaches }u\}\cdot
    \mathbf{1}\{\mathbf{D}[r,j]=b\}\cdot
    \mathbf{Y}[c,r].
\]
Prefix sums over the bin dimension give the class counts on the left and right
sides of each candidate split.  For a candidate split $(u,j,b)$, write
$N_{L,c}$ and $N_{R,c}$ for the class counts on the two sides, and
$N_L=\sum_c N_{L,c}$ and $N_R=\sum_c N_{R,c}$.  Maximizing the usual Gini
improvement is equivalent, for a fixed parent node, to maximizing
\begin{equation}\label{eq:rf-gini-score}
    \Phi_{\mathsf{rf}}(u,j,b)
    =
    \frac{\sum_{c\in[K]} N_{L,c}^2}{N_L}
    +
    \frac{\sum_{c\in[K]} N_{R,c}^2}{N_R}.
\end{equation}
The trainer selects the best valid split among the features allowed by
$\mathbf{F}^{(q)}$.  At each leaf, the output can be either the majority class
or the class-frequency vector, depending on the public trainer.

For regression random forests, the same structure applies after replacing
class-frequency histograms with count, label-sum, and label-square-sum
histograms.  The split objective then becomes the usual variance or RSS
reduction, which is the same algebraic form as the regression-tree objective in
the main construction.

\paragraph{Witness objects.}
Compared with the GBDT construction, the prover no longer needs the boosting
score tensor or the pseudo-residual tensor.  Instead, for each tree it commits
to the leaf assignment vector, the sampling multiplicities if they are not
public, the feature-availability mask if it is not public, the class-frequency
histograms, the candidate split scores, the selected feature-threshold pairs,
and the leaf outputs.

For classification random forests, the main histogram object has shape
\[
    N_{\mathrm{tot}}\times d\times B\times K.
\]
For regression random forests, the prover uses the same count/sum histogram
shape as in the GBDT protocol, with the training label replacing the
pseudo-residual target.

\paragraph{Proof shape.}
The proof again follows the same high-level skeleton.

First, the prover proves forest well-formedness and leaf assignment exactly as
in the main construction.  If the sampling masks and feature masks are derived
from a public seed, no extra proof is needed for their correctness.  If they are
committed as witness objects, the prover additionally proves that they are
consistent with the public sampling rule.

Second, the prover uses the histogram proof to certify the transformation from
sample-wise labeled records to class-frequency histograms.  This is the random
forest analogue of the histogram proof in the main protocol: instead of
aggregating pseudo-residuals, the proof aggregates class counts, optionally
weighted by bootstrap multiplicities.

Third, the prover proves the split score.  For classification, it introduces
quotient witnesses $Q_L$ and $Q_R$ and proves
\[
    Q_L\cdot N_L=\sum_{c\in[K]}N_{L,c}^2,
    \qquad
    Q_R\cdot N_R=\sum_{c\in[K]}N_{R,c}^2,
\]
together with
\[
    \Gamma=Q_L+Q_R.
\]
The validity mask enforces that $N_L$ and $N_R$ are nonzero and that the
candidate feature is allowed by the feature-availability mask.  The selected
split is then proved optimal by the same max-selection proof used in the main
construction.

Fourth, the prover proves the leaf output.  If the leaf stores a class-frequency
vector, the prover checks the corresponding quotient identities
\[
    P_{\ell,c}\cdot N_\ell=N_{\ell,c}
\]
for every class $c$.  If the leaf stores only a majority label, the prover
proves that the committed label is the class with maximum leaf count using the
same max-selection primitive.  Inference then consists of routing each input to
one leaf per tree, fetching the leaf vote or probability vector, summing across
trees, and proving the final majority or argmax decision.

\paragraph{Batching.}
The batching strategy is unchanged.  The class-frequency histograms only add one
more tensor dimension, and all tensors are flattened before commitment as in the
main protocol.  Domain-lifting batching handles the linear constraints for
histogram propagation, prefix sums, vote aggregation, and inference
aggregation.  Interleaving batching handles lookup, range, quotient, and
max-selection checks across the heterogeneous objects introduced by different
trees and nodes.

\paragraph{Completeness and soundness.}
Completeness follows because an honestly trained random forest produces leaf
assignments, histograms, split scores, selected splits, and leaf outputs
satisfying the stated relations.  Soundness follows because any accepting proof
binds the prover to committed histograms and split scores that are consistent
with the committed data, the public sampling and feature-selection rule, and
the public split objective.  Therefore, the accepted forest is consistent with a
valid execution of the specified random-forest trainer.
